\documentclass[12pt]{article}

\input{../../_assets/preambulo.tex}


\begin{document}

    % 1. Foto de fondo
    % 2. Título
    % 3. Encabezado Izquierdo
    % 4. Color de fondo
    % 5. Coord x del titulo
    % 6. Coord y del titulo
    % 7. Fecha

    
    \input{../../_assets/portada}
    \portadaExamen{ffccA4.jpg}{Variable Compleja I\\Examen XX}{Variable Compleja I. Examen XX}{MidnightBlue}{-9}{28}{2026}{David Muñoz Gómez}

    \begin{description}
        \item[Asignatura] Variable Compleja I.
        \item[Curso Académico] 2025/26.
        \item[Grado] Doble Grado en Ingeniería Informática y Matemáticas.
        \item[Grupo] Único.
        \item[Profesor] Javier Merí de la Maza.
        \item[Descripción] Prueba intermedia de Variable Compleja I.
        \item[Fecha] 13 de mayo de 2026.
        % \item[Duración] ---.
    \end{description}
    \newpage


    % ------------------------------------
    
    \begin{ejercicio}[4 puntos]
        Demostrar los siguientes enunciados:
        \begin{enumerate}[label=\alph*)]
            \item $f \in \cc{H}(\bb{C})$ con $\Re(f)$ acotada implica que $f$ es constante.
            \item Si $f \in \cc{H}(\bb{C})$ cumple que para cada $z \in \bb{C}$ alguna derivada de $f$ se anula en $z$ (esto es, para cada $z \in \bb{C}$ existe $n \in \bb{N}$ tal que $f^{(n)}(z) = 0$), entonces $f$ es un polinomio.
        \end{enumerate}
    \end{ejercicio}

    \begin{ejercicio}[3 puntos]
        Probar que la serie 
        \[ \sum_{n \geq 1} \frac{\text{sen}(n + z)}{n^z} \]
        converge absolutamente en todo punto del dominio $\Omega = \{z \in \bb{C} : \text{Re } z > 1\}$ y uniformemente en cada subconjunto compacto contenido en $\Omega$.
    \end{ejercicio}

    \begin{ejercicio}[3 puntos]
        Sean $f$ y $g \in \cc{H}(\bb{C})$ funciones no idénticamente nulas cumpliendo que $|g(z)| \leq |f(z)|$ para cada $z \in \bb{C}$. Probar que existe $\alpha \in \bb{C}$ de modo que $g(z) = \alpha f(z)$ para cada $z \in \bb{C}$.
    \end{ejercicio}

    \newpage
    \setcounter{ejercicio}{0} % Reiniciar contador de ejercicios
    \noindent

\end{document}